% GENERATED FILE. Edit canonical lessons, then run npm run generate:books.
% canonical-source-hash: fnv1a64:5e1249f7559da814
% canonical-lessons: BN-C54-banano, BN-C54-tola, BN-C54-tana, BN-C54-lukano, BN-C54-jeta

\chapter{To make, To lift, To pull, To hide, To win}
\label{ch:bn-to-make-to-lift-to-pull-to-hide-to-win}
\hlchaptermodality{54}

\begin{chapteropening}
\textbf{By the end of this chapter:} \emph{I can say to make, to lift, to pull, to hide and to win.}

Complete the last lesson of chapter 54: I can say to make, to lift, to pull, to hide and to win.
\end{chapteropening}

\section[\texorpdfstring{\bn{বানানো} (bānāno)}{বানানো (bānāno)}]{\bn{বানানো} (bānāno) — to make}

\label{lesson:BN-C54-banano}



\begin{quote}
Before the new one: say the Bengali for to fly, then the Bengali for to teach.
\end{quote}

\subsection*{You'll want to know: \bn{বানানো}}
\textbf{\bn{বানানো}} — \emph{bānāno} — \textquotedblleft{}to make\textquotedblright{}.

\begin{cousinweb}
\textbf{\bn{বানানো}} (\emph{bānāno}) is a tadbhava, a word inherited from Sanskrit and reshaped by sound change. It comes from \emph{vanati}, \textquotedblleft{}desires, makes ready\textquotedblright{} (CDIAL 11260). \emph{Bānāno} is the causing form of \emph{banā}, a verb CDIAL glosses \textquotedblleft{}to agree, do\textquotedblright{}.
\end{cousinweb}

\begin{grammarlens}[title={What you've built}]
One more word for this chapter.
\end{grammarlens}

\subsection*{Your turn}
\begin{itemize}
\raggedright
  \item \emph{Say it:} \emph{bānāno}
  \item \emph{Say it:} \emph{bānāno}, once more
  \item \emph{Say it:} say \emph{bānāno}
  \item \emph{Recall:} say the Bengali for to fly, then the Bengali for to teach, then say \emph{bānāno} again
\end{itemize}

\begin{tcolorbox}[breakable,colback=teal!4,colframe=teal!35!black,title={Before you move on}]
What does \bn{বানানো} mean? (\textquotedblleft{}To make\textquotedblright{}.) Say it once more.
\end{tcolorbox}
\section[\texorpdfstring{\bn{তোলা} (tolā)}{তোলা (tolā)}]{\bn{তোলা} (tolā) — to lift}

\label{lesson:BN-C54-tola}



\begin{quote}
Before the new one: say the Bengali for to teach, then the Bengali for to make.
\end{quote}

\subsection*{You'll want to know: \bn{তোলা}}
\textbf{\bn{তোলা}} — \emph{tolā} — \textquotedblleft{}to lift\textquotedblright{}.

\begin{cousinweb}
\textbf{\bn{তোলা}} (\emph{tolā}) is a tadbhava, a word inherited from Sanskrit and reshaped by sound change. It comes from \emph{tolayati}, \textquotedblleft{}lifts, weighs\textquotedblright{} (CDIAL 5979).
\end{cousinweb}

\begin{grammarlens}[title={What you've built}]
One more word for this chapter.
\end{grammarlens}

\subsection*{Your turn}
\begin{itemize}
\raggedright
  \item \emph{Say it:} \emph{tolā}
  \item \emph{Say it:} \emph{tolā}, once more
  \item \emph{Say it:} say \emph{tolā}
  \item \emph{Recall:} say the Bengali for to teach, then the Bengali for to make, then say \emph{tolā} again
\end{itemize}

\begin{tcolorbox}[breakable,colback=teal!4,colframe=teal!35!black,title={Before you move on}]
What does \bn{তোলা} mean? (\textquotedblleft{}To lift\textquotedblright{}.) Say it once more.
\end{tcolorbox}
\section[\texorpdfstring{\bn{টানা} (ṭānā)}{টানা (ṭānā)}]{\bn{টানা} (ṭānā) — to pull}

\label{lesson:BN-C54-tana}



\begin{quote}
Before the new one: say the Bengali for to make, then the Bengali for to lift.
\end{quote}

\subsection*{You'll want to know: \bn{টানা}}
\textbf{\bn{টানা}} — \emph{ṭānā} — \textquotedblleft{}to pull\textquotedblright{}.

\begin{cousinweb}
\textbf{\bn{টানা}} (\emph{ṭānā}) is a tadbhava, a word inherited and reshaped by sound change. It comes from \emph{*ṭān}, \textquotedblleft{}tighten\textquotedblright{}, a form CDIAL reconstructs from its descendants (CDIAL 5443). CDIAL calls its relation to Sanskrit \emph{tanoti}, \textquotedblleft{}stretches\textquotedblright{}, obscure.
\end{cousinweb}

\begin{grammarlens}[title={What you've built}]
One more word for this chapter.
\end{grammarlens}

\subsection*{Your turn}
\begin{itemize}
\raggedright
  \item \emph{Say it:} \emph{ṭānā}
  \item \emph{Say it:} \emph{ṭānā}, once more
  \item \emph{Say it:} say \emph{ṭānā}
  \item \emph{Recall:} say the Bengali for to make, then the Bengali for to lift, then say \emph{ṭānā} again
\end{itemize}

\begin{tcolorbox}[breakable,colback=teal!4,colframe=teal!35!black,title={Before you move on}]
What does \bn{টানা} mean? (\textquotedblleft{}To pull\textquotedblright{}.) Say it once more.
\end{tcolorbox}
\section[\texorpdfstring{\bn{লুকানো} (lukāno)}{লুকানো (lukāno)}]{\bn{লুকানো} (lukāno) — to hide}

\label{lesson:BN-C54-lukano}



\begin{quote}
Before the new one: say the Bengali for to lift, then the Bengali for to pull.
\end{quote}

\subsection*{You'll want to know: \bn{লুকানো}}
\textbf{\bn{লুকানো}} — \emph{lukāno} — \textquotedblleft{}to hide\textquotedblright{}.

\begin{cousinweb}
\textbf{\bn{লুকানো}} (\emph{lukāno}) is a tadbhava, a word inherited from Sanskrit and reshaped by sound change. It comes from \emph{lupyate}, \textquotedblleft{}is lost, is removed\textquotedblright{} (CDIAL 11083). \emph{Lukāno} is the causing form of \emph{lukā}, \textquotedblleft{}to hide\textquotedblright{}, the verb CDIAL lists.
\end{cousinweb}

\begin{grammarlens}[title={What you've built}]
One more word for this chapter.
\end{grammarlens}

\subsection*{Your turn}
\begin{itemize}
\raggedright
  \item \emph{Say it:} \emph{lukāno}
  \item \emph{Say it:} \emph{lukāno}, once more
  \item \emph{Say it:} say \emph{lukāno}
  \item \emph{Recall:} say the Bengali for to lift, then the Bengali for to pull, then say \emph{lukāno} again
\end{itemize}

\begin{tcolorbox}[breakable,colback=teal!4,colframe=teal!35!black,title={Before you move on}]
What does \bn{লুকানো} mean? (\textquotedblleft{}To hide\textquotedblright{}.) Say it once more.
\end{tcolorbox}
\section[\texorpdfstring{\bn{জেতা} (jetā)}{জেতা (jetā)}]{\bn{জেতা} (jetā) — to win}

\label{lesson:BN-C54-jeta}



\begin{quote}
Before the new one: say the Bengali for to pull, then the Bengali for to hide.
\end{quote}

\subsection*{You'll want to know: \bn{জেতা}}
\textbf{\bn{জেতা}} — \emph{jetā} — \textquotedblleft{}to win\textquotedblright{}.

\begin{cousinweb}
\textbf{\bn{জেতা}} (\emph{jetā}) is a tadbhava, a word inherited from Sanskrit and reshaped by sound change. It comes from \emph{jita}, \textquotedblleft{}won, conquered\textquotedblright{} (CDIAL 5224). CDIAL writes the Bengali verb \emph{jitā}.
\end{cousinweb}

\begin{grammarlens}[title={What you've built}]
One more word for this chapter.
\end{grammarlens}

\subsection*{Your turn}
\begin{itemize}
\raggedright
  \item \emph{Say it:} \emph{jetā}
  \item \emph{Say it:} \emph{jetā}, once more
  \item \emph{Say it:} say \emph{jetā}
  \item \emph{Recall:} say the Bengali for to pull, then the Bengali for to hide, then say \emph{jetā} again
\end{itemize}

\begin{tcolorbox}[breakable,colback=teal!4,colframe=teal!35!black,title={Before you move on}]
What does \bn{জেতা} mean? (\textquotedblleft{}To win\textquotedblright{}.) Say it once more.
\end{tcolorbox}
